\documentclass[11pt,a4paper]{article}

\usepackage[utf8]{inputenc}
\usepackage[T1]{fontenc}
\usepackage{lmodern}
\usepackage{microtype}
\usepackage[margin=1in]{geometry}
\usepackage{booktabs}
\usepackage{longtable}
\usepackage{array}
\usepackage{enumitem}
\usepackage{hyperref}
\usepackage{url}
\usepackage{xcolor}
\usepackage{listings}
\usepackage{fvextra}

\hypersetup{
  colorlinks=true,
  linkcolor=black,
  citecolor=blue,
  urlcolor=blue,
  pdftitle={MentisDB: Durable Semantic Memory for Software Agents},
  pdfauthor={Angel Leon}
}

\lstdefinelanguage{Rust}{
  keywords={fn, let, mut, struct, enum, impl, for, in, if, else, return, match, use,
           pub, crate, mod, self, Some, None, Ok, Err, as, where, dyn, trait, type,
           const, static, ref, move, async, await, break, continue, loop, while,
           true, false, Vec, HashMap, HashSet, Option, Result, String, Box,
           VecDeque, Uuid, DateTime, Utc, usize, u64, u32, f64, bool, str},
  morecomment=[l]{//},
  morecomment=[s]{/*}{*/},
  morestring=[b]",
  sensitive=true,
}

\lstset{
  language=Rust,
  basicstyle=\ttfamily\small,
  keywordstyle=\color{blue!70!black}\bfseries,
  commentstyle=\color{green!50!black}\itshape,
  stringstyle=\color{red!60!black},
  numberstyle=\tiny\color{gray},
  showstringspaces=false,
  breaklines=true,
  breakatwhitespace=true,
  frame=single,
  framerule=0.4pt,
  rulecolor=\color{gray!50},
  backgroundcolor=\color{gray!5},
  tabsize=2,
  columns=flexible,
  literate={→}{$\to$}{1}
           {≥}{$\geq$}{1}
           {≤}{$\leq$}{1}
           {≠}{$\neq$}{1},
  xleftmargin=0.5em,
  xrightmargin=0.5em,
}

\title{\textbf{MentisDB}\\[0.3em]
\large Durable Semantic Memory for Software Agents}
\author{Angel Leon\\[0.3em]
\small Version 0.10.8.53 \quad \textemdash \quad September 16, 2026}
\date{}

\begin{document}
\maketitle

\begin{abstract}
Most agent frameworks treat long-term memory as an afterthought — ad-hoc prompt stuffing,
unstructured Markdown files, or opaque session state that's non-portable and easily lost.
MentisDB is a durable, semantically typed memory engine built in Rust that treats agent
memory as an append-only, hash-chained ledger of structured records called \emph{thoughts}.
Every thought is cryptographically linked to its predecessor via SHA-256, making the chain
tamper-evident. On top of this ledger, MentisDB provides a retrieval pipeline that combines
BM25 lexical scoring, vector similarity, typed graph expansion, implicit edge inference,
and Reciprocal Rank Fusion (RRF) — all without requiring an external database or LLM service
in the core path. On standard long-term memory benchmarks, MentisDB achieves 72.6\% Recall@10
on LoCoMo-10P and 66.8\% Recall@5 on LongMemEval.
\end{abstract}

\tableofcontents
\newpage

% ─────────────────────────────────────────────────────────────────────────────
\section{System Overview}

\subsection{The Problem}

LLM agents need memory that is:
\begin{itemize}
  \item \textbf{Durable} — survives process restarts and context window resets
  \item \textbf{Queryable} — not just a flat file, but searchable by text, semantics, and graph relationships
  \item \textbf{Tamper-evident} — you can detect if someone modified the memory log
  \item \textbf{Portable} — works with Claude Code, Codex, Copilot, Cursor, or any MCP client
  \item \textbf{Attributable} — every record carries the agent ID that wrote it
\end{itemize}

Existing solutions fail one or more of these. Markdown files aren't queryable. Provider
session state isn't portable. Vector databases don't have integrity guarantees. LLM-based
memory extraction adds latency and cost to every write.

\subsection{The Approach}

MentisDB models agent memory as a hash-chained append-only log — similar to a blockchain,
but without consensus or distributed systems complexity. Each entry (a ``thought'') is a
typed, structured record with semantic metadata, signed by the agent that produced it.

Retrieval is a multi-signal pipeline: lexical search (BM25), vector similarity (cosine
over ONNX embeddings), graph traversal (typed edges between thoughts), and rank fusion
(RRF). All of this runs locally — no network calls, no LLM API bills for the core path.

\subsection{Architecture at a Glance}

The daemon is a single static binary with no external dependencies. Storage is embedded
files on disk. No database server, no Redis, no Neo4j. The MCP server runs on port 9471,
REST on 9472, and the HTTPS dashboard on 9473. All storage lives under
\texttt{\textasciitilde/.cloudllm/mentisdb/}.

% ─────────────────────────────────────────────────────────────────────────────
\section{The Thought Data Model}

\subsection{The Thought Record}

The core unit of memory is a \emph{thought} — a structured, typed, hash-chained record:

\begin{lstlisting}
struct Thought {
    schema_version: u32,

    // Identity (assigned by the chain on commit)
    thought_id: Uuid,
    index: u64,
    timestamp: DateTime<Utc>,

    // Attribution
    agent_id: String,
    signing_key_id: Option<String>,
    thought_signature: Option<Vec<u8>>,

    // Semantics
    thought_type: ThoughtType,
    thought_role: ThoughtRole,
    content: String,
    tags: Vec<String>,
    concepts: Vec<String>,

    // Scoring metadata
    confidence: Option<f64>,
    importance: Option<f64>,

    // Relationships
    refs: Vec<u64>,
    relations: Vec<ThoughtRelation>,

    // Chain integrity (assigned on commit)
    prev_hash: String,
    hash: String,
}
\end{lstlisting}

The key design decision: the agent authors the \emph{content} fields but cannot forge the
\emph{chain} fields (id, index, timestamp, hashes). Those are assigned by the chain on commit.

\subsection{The Hash Chain}

Every thought is linked to its predecessor via SHA-256, similar to a git history but for
memory records:

\begin{lstlisting}
fn compute_thought_hash(thought: &Thought) -> String {
    // Serialize everything EXCEPT the hash field itself
    let canonical_bytes = bincode::serialize(&CanonicalThought {
        schema_version: thought.schema_version,
        thought_id: thought.thought_id,
        index: thought.index,
        // ... all fields except `hash`
        prev_hash: thought.prev_hash.clone(),
    });
    sha256(&canonical_bytes)
}
\end{lstlisting}

The integrity invariant: \texttt{thought.hash == sha256(canonical\_serialize(thought))}
and \texttt{thought.prev\_hash == previous\_thought.hash}.

If someone modifies any field, the hash breaks. If they recompute the hash, the next
thought's \texttt{prev\_hash} no longer matches — tamper cascades forward through the
entire chain. Detection is O(1): check that each record's \texttt{prev\_hash} matches the
previous record's \texttt{hash}.

For stronger provenance, individual thoughts can be Ed25519-signed by the producing agent.

\subsection{Typed Relations}

Thoughts link to each other via 12 typed edge kinds:

\begin{lstlisting}
enum ThoughtRelationKind {
    References, Summarizes, Corrects, Invalidates,
    CausedBy, Supports, Contradicts, DerivedFrom,
    ContinuesFrom, BranchesFrom, RelatedTo, Supersedes,
}
\end{lstlisting}

Each relation carries an optional validity interval (\texttt{valid\_at} /
\texttt{invalid\_at}), enabling point-in-time queries: ``what did the agent know at
timestamp T?''

\subsection{The Invalidation Set}

When a thought is superseded, corrected, or invalidated, we don't delete it — we mark it.
At chain open time, we precompute a set of all invalidated thought IDs for O(1) skipping
during retrieval:

\begin{lstlisting}
fn build_invalidation_set(chain: &[Thought]) -> HashSet<Uuid> {
    let mut invalidated = HashSet::new();
    for thought in chain {
        for relation in &thought.relations {
            match relation.kind {
                Supersedes | Corrects | Invalidates => {
                    invalidated.insert(relation.target_thought_id);
                }
                _ => {}
            }
        }
    }
    invalidated
}
\end{lstlisting}

\subsection{Agent Registry}

Thoughts carry only a stable \texttt{agent\_id} string; display names, owners, descriptions,
aliases, public keys, and thought counters live in a per-chain agent registry, resolved at
read time.

% ─────────────────────────────────────────────────────────────────────────────
\section{Semantic Typing}

\subsection{ThoughtType}

Every thought carries a semantic type from a 30-variant enum, organized into seven
categories: user/relationship, observation, error/correction, planning, exploration,
execution, and state. Examples include \texttt{Decision}, \texttt{Mistake},
\texttt{Insight}, \texttt{Constraint}, \texttt{Checkpoint}, and \texttt{Summary}.

\subsection{ThoughtRole}

Orthogonal to type, every thought has a role: \texttt{Memory}, \texttt{WorkingMemory},
\texttt{Summary}, \texttt{Compression}, \texttt{Checkpoint}, \texttt{Handoff},
\texttt{Audit}, \texttt{Retrospective}, or \texttt{Dream}. The combination of type
$\times$ role gives 270 distinguishable semantic positions.

% ─────────────────────────────────────────────────────────────────────────────
\section{Storage Layer}

\subsection{Storage Adapter}

Persistence is abstracted behind a trait, allowing the chain semantics to stay the same
regardless of backend:

\begin{lstlisting}
trait StorageAdapter {
    fn load_thoughts(&self) -> io::Result<Vec<Thought>>;
    fn append_thought(&self, thought: &Thought) -> io::Result<()>;
    fn flush(&self) -> io::Result<()>;
    fn set_auto_flush(&self, auto: bool);
}
\end{lstlisting}

\subsection{Binary Storage}

The default backend writes each thought as a length-prefixed bincode record with file
extension \texttt{.tcbin}. This format is 40–60\% smaller than JSON and deserializes faster.

Two durability modes: \textbf{Strict} (each append blocks until flushed; 2ms group-commit
window) and \textbf{Buffered} (flushes every 16 records; up to 15 may be lost on hard crash).

% ─────────────────────────────────────────────────────────────────────────────
\section{Schema Evolution}

MentisDB has a linear schema version history (V0 through V3). Each migration is an
idempotent transformation — running it twice produces the same result as once. Migrations
compose: a V0 chain is upgraded by applying V0$\to$V1, then V1$\to$V2, then V2$\to$V3.

A critical constraint: because bincode encodes enum variants by integer tag, new variants
must be appended to the end of an enum. Reordering or inserting mid-enum would silently
corrupt persisted data.

% ─────────────────────────────────────────────────────────────────────────────
\section{Retrieval Pipeline}

\subsection{Baseline Filter Search}

The baseline path narrows candidates by indexed fields (type, role, agent\_id, tags,
concepts) and applies a case-insensitive substring predicate. Results return in append
order. No BM25, no vectors, no graph.

\subsection{Ranked Search — Backend Selection}

Ranked search selects a backend based on available signals: Lexical (BM25 only), Hybrid
(BM25 + vector fusion), LexicalGraph, HybridGraph (full pipeline), or Heuristic (when no
text is provided).

\subsection{BM25 with Per-Field DF Gating}

The lexical score uses BM25 applied across five fields. Each field has its own weight and
document-frequency gate:

\begin{lstlisting}
fn bm25_field_score(
    term_freq: f64,
    doc_field_len: f64,
    avg_field_len: f64,
    doc_freq: usize,
    corpus_size: usize,
) -> f64 {
    let k1 = 1.2;
    let b = 0.75;
    let idf = ((corpus_size as f64 - doc_freq as f64 + 0.5)
        / (doc_freq as f64 + 0.5) + 1.0).ln();
    let tf_norm = (term_freq * (k1 + 1.0))
        / (term_freq + k1 * (1.0 - b + b * (doc_field_len / avg_field_len)));
    idf * tf_norm
}
\end{lstlisting}

Per-field DF gating prevents common terms from dominating. If a term appears in more than
30\% of documents in a field, it's filtered out of that field's scoring. Each field has its
own cutoff: content (30\%, weight 1.0), tags (30\%, weight 1.6), concepts (30\%, weight 1.4),
agent\_id (70\%, weight 1.5), agent\_registry (60\%, weight 1.1).

Text normalization applies Porter stemming plus an irregular verb lemma table ($\sim$170
entries: ``went'' $\to$ ``go'', ``saw'' $\to$ ``see'').

\subsection{Smooth Vector-Lexical Fusion}

When a vector sidecar is available, cosine similarity provides a semantic score. The fusion
function combines lexical and vector signals:

\begin{lstlisting}
fn fuse_scores(lexical_score: f64, vector_similarity: f64) -> f64 {
    let alpha = 35.0;
    let beta = 3.0;
    vector_similarity * (1.0 + alpha * (-lexical_score / beta).exp())
}
\end{lstlisting}

When lexical score is 0 (no text overlap), the vector signal gets amplified $\sim$36$\times$.
As lexical score increases, amplification decays exponentially, so strong lexical matches
aren't drowned out.

\subsection{Graph-Aware Expansion}

Thoughts linked by typed relations form a directed graph. Ranked search expands outward
from seed hits via bounded BFS, traversing both explicit typed edges and implicit
cosine-inferred edges:

\begin{lstlisting}
fn expand_graph(
    seeds: &[&Thought],
    adjacency: &AdjacencyIndex,
    implicit_edges: &ImplicitEdgeOverlay,
    max_depth: usize,
    max_visits: usize,
    mode: TraversalMode,
) -> Vec<GraphHit> {
    let mut visited = HashSet::new();
    let mut queue = VecDeque::new();
    let mut results = Vec::new();

    for seed in seeds {
        queue.push_back((seed, 0));
    }

    while let Some((thought, depth)) = queue.pop_front() {
        if depth > max_depth || results.len() >= max_visits {
            break;
        }
        if !visited.insert(thought.thought_id) {
            continue;
        }
        results.push(GraphHit {
            thought_id: thought.thought_id,
            depth,
            graph_score: 1.0 / (depth as f64),
        });
        for neighbor in adjacency.neighbors(thought, mode) {
            queue.push_back((neighbor, depth + 1));
        }
        for neighbor in implicit_edges.neighbors(thought) {
            queue.push_back((neighbor, depth + 1));
        }
    }
    results
}
\end{lstlisting}

Each typed relation carries a weight: \texttt{ContinuesFrom} (0.60) is stronger than
\texttt{References} (0.06). Graph proximity decays with depth: depth 1 gets 1.0, depth 2
gets 0.5, depth 3 gets 0.33.

\subsection{Session Cohesion}

When a seed thought has a moderate lexical score (not strong enough to stand alone),
thoughts adjacent in append order get a proximity boost, surfacing ``evidence turns'' that
share context but not vocabulary:

\begin{lstlisting}
fn cohesion_score(
    candidate_index: u64,
    seed_index: u64,
    seed_lexical_score: f64,
) -> f64 {
    let radius = 8;
    let boost = 0.8;
    if seed_lexical_score < 3.0 || seed_lexical_score >= 5.0 {
        return 0.0;
    }
    let distance = (candidate_index as f64 - seed_index as f64).abs();
    if distance > radius as f64 {
        return 0.0;
    }
    boost * (1.0 - distance / radius as f64)
}
\end{lstlisting}

\subsection{Implicit Edge Overlay}

Most agents append thoughts without authoring explicit relation edges. The implicit edge
overlay derives \texttt{RelatedTo} edges automatically from vector cosine similarity:

\begin{lstlisting}
fn build_implicit_edges(
    vectors: &HashMap<Uuid, Vec<f32>>,
    threshold: f64,      // default 0.85
    max_neighbors: usize, // default 5
) -> HashMap<Uuid, Vec<ImplicitNeighbor>> {
    let mut overlay = HashMap::new();
    let ids: Vec<_> = vectors.keys().collect();

    for &source_id in &ids {
        let source_vec = &vectors[source_id];
        let mut neighbors: Vec<(Uuid, f64)> = Vec::new();

        for &target_id in &ids {
            if source_id == target_id { continue; }
            let sim = cosine_similarity(source_vec, &vectors[target_id]);
            if sim >= threshold {
                neighbors.push((*target_id, sim));
            }
        }
        neighbors.sort_by(|a, b| b.1.partial_cmp(&a.1).unwrap());
        neighbors.truncate(max_neighbors);
        overlay.insert(*source_id, neighbors.into_iter()
            .map(|(id, score)| ImplicitNeighbor { id, score })
            .collect());
    }
    overlay
}
\end{lstlisting}

The full build is O(N$^2$), done once at chain open for N $\leq$ 50,000. On each new
append, an O(N) incremental update computes cosines between the new vector and all existing
entries.

\subsection{HNSW Approximate Vector Backend}

For vector sidecars above 50,000 vectors, MentisDB switches from exact linear-scan cosine
to an HNSW graph. Both backends implement the same \texttt{VectorSearchBackend} trait, so
the retrieval pipeline is unchanged. At 10k vectors, HNSW delivers a 10.9$\times$ query
speedup with 91.4\% recall. At 50k vectors, it achieves 100\% recall at 4.23ms per query.

\subsection{Reciprocal Rank Fusion}

When reranking is enabled, the top K candidates are re-ranked using RRF — a simple method
that combines multiple ranked lists:

\begin{lstlisting}
fn rrf_score(
    candidate: &Thought,
    ranked_lists: &[Vec<Uuid>],
    k: u64,  // damping constant, default 60
) -> f64 {
    let mut total = 0.0;
    for list in ranked_lists {
        if let Some(pos) = list.iter().position(|id| *id == candidate.thought_id) {
            let rank = (pos + 1) as u64;
            total += 1.0 / (k + rank) as f64;
        }
    }
    total
}
\end{lstlisting}

MentisDB produces three single-signal rankings — lexical-only, vector-only, and
graph-only — and fuses them via RRF. Pure arithmetic: no LLM, no network round trip.

\subsection{Importance Weighting}

User-originated thoughts (importance $\sim$0.8) outrank verbose assistant responses
(importance $\sim$0.2) via a differential boost:

\begin{lstlisting}
fn importance_boost(base_score: f64, importance: f64) -> f64 {
    base_score * (importance - 0.5) * 0.3
}
\end{lstlisting}

\subsection{Decomposed Scores}

Each ranked hit exposes its full score breakdown (lexical, vector, graph, relation boost,
seed score, importance, confidence, recency, cohesion, total) plus matched terms and match
sources, preserving auditability.

% ─────────────────────────────────────────────────────────────────────────────
\section{Deduplication}

When a new thought is appended, MentisDB checks the last N thoughts (default 64) for
near-duplicates using Jaccard similarity over normalized token sets:

\begin{lstlisting}
fn check_dedup(
    new_thought: &Thought,
    recent_thoughts: &[&Thought],
    threshold: f64,  // default 0.85
) -> Option<Uuid> {
    let new_tokens = normalize_tokens(&new_thought.content);
    if new_tokens.is_empty() { return None; }

    let mut best_match = None;
    let mut best_sim = 0.0;

    for prior in recent_thoughts {
        let prior_tokens = normalize_tokens(&prior.content);
        let sim = jaccard(&new_tokens, &prior_tokens);
        if sim > best_sim {
            best_sim = sim;
            best_match = Some(prior.thought_id);
        }
    }

    if best_sim >= threshold { best_match } else { None }
}

fn jaccard(a: &HashSet<String>, b: &HashSet<String>) -> f64 {
    let intersection = a.intersection(b).count() as f64;
    let union = a.union(b).count() as f64;
    if union == 0.0 { 0.0 } else { intersection / union }
}
\end{lstlisting}

If a duplicate is found, a \texttt{Supersedes} edge is auto-emitted and the superseded
thought's ID is added to the invalidation set. The original is retained for audit.

% ─────────────────────────────────────────────────────────────────────────────
\section{Operational Surfaces}

\subsection{CLI}

The \texttt{mentisdb} binary exposes subcommands that RPC over REST to a running daemon:
\texttt{add}, \texttt{search}, \texttt{agents}. Client commands use synchronous HTTP to
avoid an async runtime dependency.

\subsection{MCP Server}

MentisDB exposes a streamable HTTP MCP endpoint (port 9471) with 35+ tools covering
bootstrap, append, search, read, export/import, agent registry, chain management, and a
skill registry.

\subsection{Skill Registry}

A git-like immutable version store for agent instruction bundles. Uploads to an existing
skill ID create new immutable versions: the first is full content, subsequent versions are
unified diff patches. Skills may be deprecated or revoked without losing history.
Permanent delete removes the skill and all versions so the identifier can be reused.
Agents with registered Ed25519 keys must sign uploads.

\subsection{Bootstrap Protocol}

Modern MCP clients bootstrap from a handshake: \texttt{initialize} directs the agent to
read the core skill, \texttt{mentisdb\_bootstrap} opens or creates the chain (writing a
genesis checkpoint if empty), and \texttt{mentisdb\_recent\_context} loads prior state.

\subsection{REST API}

A versioned REST API (\texttt{/v1/...}) mirrors the MCP tools for non-MCP clients, with
bearer token authorization supporting chain-scoped tokens.

% ─────────────────────────────────────────────────────────────────────────────
\section{Empirical Evaluation}

\subsection{Benchmarks}

\begin{longtable}{llcccc}
\toprule
Benchmark & Metric & v0.8.1 & v0.8.5 & v0.8.9 \\
\midrule
LoCoMo-2P & R@10 & \textbf{88.7\%} & — & — \\
LoCoMo-2P single-hop & R@10 & 90.7\% & — & — \\
LoCoMo-10P (1977 queries) & R@10 & 74.2\% & \textbf{74.6\%} & \textbf{71.9\%} \\
LoCoMo-10P single-hop & R@10 & — & 79.0\% & 75.8\% \\
LoCoMo-10P multi-hop & R@10 & — & 58.4\% & 57.4\% \\
LoCoMo-10P & R@20 & — & — & 79.1\% \\
LongMemEval (fresh chain) & R@5 & 67.6\% & — & \textbf{66.8\%} \\
LongMemEval (fresh chain) & R@10 & 73.2\% & — & \textbf{72.2\%} \\
LongMemEval (fresh chain) & R@20 & — & — & 78.0\% \\
\bottomrule
\end{longtable}

\subsection{Scoring Evolution}

\begin{longtable}{llcc}
\toprule
Version & Change & LongMemEval R@5 & LoCoMo-10P R@10 \\
\midrule
0.8.0 baseline & — & 57.2\% & — \\
0.8.0 + Porter stemming & token normalization & 61.6\% & — \\
0.8.0 + tiered fusion + importance & vector/lexical balance & 65.0\% & — \\
0.8.1 + cohesion + smooth fusion + DF cutoff & retrieval quality & 67.6\% & 74.2\% \\
0.8.5 + cohesion tuning + edge weights + fastembed & session/graph boost & — & 74.6\% \\
0.8.9 + irregular lemmas + webhooks & lemma expansion + events & 66.8\% & 71.9\% \\
0.9.8 + implicit edge overlay + in-memory vector index & graph density + latency & 66.8\% & — \\
0.9.9 + thesaurus auto-by-default + full embeddings & automatic synonym expansion & 66.8\% & \textbf{72.6\%} \\
\bottomrule
\end{longtable}

\subsection{Near-Miss Analysis (LoCoMo-10P, v0.8.5)}

Of 503 misses (gold answer absent from top-10):

\begin{tabular}{lrl}
\toprule
Bucket & Count & Interpretation \\
\midrule
R@20 hit & 130 (25.8\%) & close ranking error \\
R@50 hit & 285 (56.7\%) & moderate signal gap \\
R > 50 & 218 (43.3\%) & lexical gap \\
\bottomrule
\end{tabular}

The 43.3\% figure represents a hard ceiling for BM25-only retrieval. The 25.8\%
ranking-error bucket is the target for graph-density improvements.

\subsection{Near-Miss Analysis (LongMemEval, v0.9.8)}

R@5 = 66.8\%, R@10 $\approx$ 72.0\%, R@20 $\approx$ 77.8\%. Of 166 misses:

\begin{tabular}{lrl}
\toprule
Bucket & Count & Interpretation \\
\midrule
R@10 hit & 27 (16.3\%) & close ranking error \\
R@20 hit & 57 (34.3\%) & moderate signal gap \\
not in R@20 & 109 (65.7\%) & lexical gap \\
\bottomrule
\end{tabular}

65.7\% of misses are not recoverable at R@20 regardless of graph topology — evidence
thoughts share no vocabulary with the query.

\subsection{HNSW Performance}

\begin{tabular}{lllll}
\toprule
Corpus & Backend & Query Latency & Recall@10 & Speedup \\
\midrule
10k $\times$ 128d & Exact (linear scan) & 95.1 ms & 100.0\% & — \\
10k $\times$ 128d & \textbf{HNSW} & \textbf{8.75 ms} & 91.4\% & \textbf{10.9$\times$} \\
50k $\times$ 128d & \textbf{HNSW} & \textbf{4.23 ms/query} & \textbf{100.0\%} & — \\
\bottomrule
\end{tabular}

\subsection{Micro-Benchmarks}

Criterion benchmarks span append throughput, baseline search, ranked retrieval, skill
registry lifecycle, and HTTP concurrency. A \texttt{DashMap}-based concurrent chain lookup
delivers 750–930 read req/s at 10k concurrent tasks.

% ─────────────────────────────────────────────────────────────────────────────
\section{Related Work and Positioning}

\begin{longtable}{lcccc}
\toprule
Feature & \textbf{MentisDB} & Mem0 & Graphiti / Zep & Letta / MemGPT \\
\midrule
Language & Rust & Python & Python & Python / TS \\
Storage & embedded file & external DB & Neo4j / FalkorDB & external DB \\
LLM required for core & \textbf{No} & Yes & Yes & Yes \\
Cryptographic integrity & \textbf{SHA-256 hash chain} & — & — & — \\
Hybrid retrieval & BM25 + vector + graph & vector + keyword & semantic + keyword + graph & — \\
Approximate NN (HNSW) & \textbf{Yes (default)} & — & — & — \\
Implicit graph inference & \textbf{cosine-inferred edges} & — & LLM-extracted & — \\
Temporal facts & valid\_at / invalid\_at & update-only & valid\_at / invalid\_at & — \\
Deduplication & \textbf{Jaccard + Supersedes} & LLM-based & merge & — \\
Agent registry & Yes & — & — & Yes \\
MCP server & \textbf{Built-in} & — & Yes & — \\
\bottomrule
\end{longtable}

% ─────────────────────────────────────────────────────────────────────────────
\section{Limitations and Future Work}

\subsection{Limitations}

\begin{itemize}
  \item \textbf{Ceiling of sparse retrieval.} 43.3\% of LoCoMo-10P misses and 65.7\% of
    LongMemEval misses are irreducible lexical gaps.
  \item \textbf{Implicit edge quality depends on the encoder.} MiniLM (384-dim) at
    threshold 0.85 produces sparse overlays for topically diverse conversations.
  \item \textbf{Local-only integrity.} The hash chain provides tamper evidence, not
    Byzantine fault tolerance.
  \item \textbf{Schema churn discipline.} Bincode enum tagging requires append-only
    variant ordering.
\end{itemize}

\subsection{Future Work}

\begin{itemize}
  \item LoCoMo multi-hop validation of implicit edge overlay
  \item Threshold sweep for cosine threshold and K
  \item Per-chain entity/relation ontologies
  \item Cross-chain federated retrieval
  \item Optional LLM-extracted memories as a layered transform
  \item Self-improving skill registry with signed provenance
\end{itemize}

\subsection{Conclusion}

MentisDB formalizes agent memory as an append-only, hash-chained ledger of semantically
typed thoughts, coupled with a composable retrieval pipeline — BM25 with per-field DF
gating, smooth vector-lexical fusion, bounded graph expansion with cosine-inferred implicit
edges, RRF, session cohesion, and Jaccard deduplication — in a single embedded Rust crate.
Empirical results demonstrate competitive retrieval quality without external databases or
LLM services for the core path.

% ─────────────────────────────────────────────────────────────────────────────
\section*{References}
\addcontentsline{toc}{section}{References}

\begin{enumerate}
  \item Robertson, S. and Zaragoza, H. \emph{The Probabilistic Relevance Framework: BM25 and Beyond}. Foundations and Trends in IR, 3(4), 2009.
  \item Cormack, G. V., Clarke, C. L. A., and Büttcher, S. \emph{Reciprocal Rank Fusion outperforms Condorcet and individual Rank Learning Methods}. SIGIR, 2009.
  \item Porter, M. F. \emph{An algorithm for suffix stripping}. Program, 14(3), 1980.
  \item Jaccard, P. \emph{Étude comparative de la distribution florale dans une portion des Alpes et des Jura}. Bull. Soc. Vaud. Sci. Nat., 1901.
  \item Bernstein, D. J. et al. \emph{High-speed high-security signatures}. J. Cryptographic Engineering, 2012.
  \item FIPS PUB 180-4. \emph{Secure Hash Standard (SHS)}. NIST, 2015.
  \item Maharana, A. et al. \emph{Evaluating Very Long-Term Conversational Memory of LLM Agents (LoCoMo)}. 2024.
  \item Wu, D. et al. \emph{LongMemEval: Benchmarking Chat Assistants on Long-Term Interactive Memory}. 2024.
  \item Anthropic. \emph{Model Context Protocol (MCP) Specification}. 2024.
\end{enumerate}

\end{document}
